\documentclass[../DoAn.tex]{subfiles}
\begin{document}
\section{Đặt vấn đề}
\label{section:1.1}

  Trong những năm gần đây, thương mại điện tử tại Việt Nam có tốc độ tăng trưởng
  mạnh mẽ và ổn định. Theo báo cáo của Hiệp hội Thương mại điện tử Việt Nam (VECOM),
  doanh thu thương mại điện tử bán lẻ năm 2023 đạt hơn 20 tỷ USD, trong đó lĩnh vực
  thời trang và phụ kiện là một trong những ngành hàng có tỷ lệ tăng trưởng cao nhất.
  Xu hướng này phản ánh sự thay đổi rõ rệt trong hành vi tiêu dùng: người mua hàng
  ngày càng ưu tiên mua sắm trực tuyến thay vì đến trực tiếp cửa hàng.

  Tuy nhiên, phần lớn doanh nghiệp bán lẻ vừa và nhỏ hiện nay vẫn phụ thuộc hoàn
  toàn vào các sàn thương mại điện tử trung gian như Shopee, Lazada hay TikTok Shop
  để tiếp cận khách hàng. Mô hình này tồn tại nhiều bất cập đáng kể. Về mặt chi phí,
  doanh nghiệp phải chịu phí hoa hồng từ 5\% đến 20\% trên mỗi đơn hàng, cùng với
  các khoản phí quảng cáo để duy trì độ hiển thị sản phẩm. Về mặt dữ liệu, toàn bộ
  thông tin khách hàng thuộc quyền kiểm soát của sàn, khiến doanh nghiệp không thể
  chủ động xây dựng mối quan hệ lâu dài với người mua. Về mặt vận hành, doanh nghiệp
  bị giới hạn trong các quy trình cứng nhắc của sàn, không thể tùy chỉnh nghiệp vụ
  đặc thù như cấu hình phí vận chuyển theo từng địa bàn, thiết lập quy trình hoàn trả
  hàng riêng, hay phân quyền linh hoạt cho từng nhóm nhân viên.

  Bài toán đặt ra là cần xây dựng một nền tảng bán hàng trực tuyến độc lập, giúp
  doanh nghiệp bán lẻ thời trang vừa và nhỏ chủ động quản lý toàn bộ vòng đời đơn
  hàng từ lúc khách hàng đặt mua đến khi hoàn tất giao dịch hoặc hoàn trả hàng hóa,
  đồng thời kiểm soát được dữ liệu khách hàng và tối ưu hóa chi phí vận hành. Nếu
  được giải quyết, bài toán này mang lại lợi ích trực tiếp cho ba nhóm đối tượng:
  chủ doanh nghiệp có thể giảm chi phí trung gian và nắm bắt hành vi khách hàng;
  nhân viên vận hành được hỗ trợ bởi công cụ quản lý rõ ràng theo vai trò; khách
  hàng được trải nghiệm mua sắm nhất quán và minh bạch hơn. Giải pháp này không
  chỉ áp dụng cho ngành thời trang mà còn có thể mở rộng sang các lĩnh vực bán lẻ
  trực tuyến khác có đặc thù nghiệp vụ tương tự.

  \section{Mục tiêu và phạm vi đề tài}
  \label{section:1.2}

  Hiện tại, các doanh nghiệp bán lẻ trực tuyến có thể lựa chọn giữa hai nhóm giải
  pháp chính. Nhóm thứ nhất là các nền tảng SaaS thương mại điện tử như Haravan,
  Sapo hay Shopify, vốn cung cấp sẵn bộ tính năng chuẩn và cho phép triển khai nhanh.
  Nhóm thứ hai là xây dựng hệ thống riêng theo yêu cầu đặc thù của từng doanh nghiệp.
  Các nền tảng SaaS có ưu điểm triển khai nhanh và chi phí ban đầu thấp, nhưng tồn
  tại các hạn chế đáng kể: chi phí thuê bao định kỳ tích lũy theo thời gian, khả năng
  tùy chỉnh nghiệp vụ bị giới hạn bởi nền tảng, và dữ liệu vẫn được lưu trữ trên
  hệ thống của bên thứ ba. Đặc biệt, các nền tảng này thường thiếu các tính năng
  chuyên biệt như quản lý hoàn trả hàng có vòng đời trạng thái rõ ràng, cấu hình
  phí vận chuyển linh hoạt đến cấp xã/phường, hay hệ thống phân quyền nhân viên
  theo vai trò kết hợp với nhật ký tác động hệ thống và thống kê.

  Trên cơ sở phân tích các hạn chế trên, đồ án này đặt mục tiêu xây dựng một ứng
  dụng web thương mại điện tử hoàn chỉnh dành cho cửa hàng thời trang và phụ kiện
  quy mô vừa và nhỏ. Hệ thống hướng đến việc giải quyết các hạn chế đã xác định
  thông qua các chức năng chính sau: quản lý sản phẩmvà danh mục; quản lý giỏ hàng, đặt hàng và thanh toán trực tuyến qua
  cổng VNPay và hình thức thanh toán khi nhận hàng (COD); quản lý vòng đời đơn hàng
  và quy trình hoàn trả hàng có nhiều trạng thái; cấu hình phí vận chuyển theo
  quận/huyện và xã/phường; quản lý mã giảm giá, danh sách yêu thích và blog;
  phân quyền nhân viên theo vai trò kết hợp với nhật ký kiểm toán hành động;
  hỗ trợ khách hàng qua chatbot theo thời gian thực. Phạm vi của đồ án tập trung
  vào việc xây dựng hệ thống back-office dành cho quản trị viên và nhân viên,
  cùng với giao diện mua sắm dành cho khách hàng cuối, không bao gồm ứng dụng di động.

  \section{Định hướng giải pháp}
  \label{section:1.3}

  Để xây dựng hệ thống với các yêu cầu đã xác định ở phần \ref{section:1.2}, đồ
  án lựa chọn hướng tiếp cận phát triển ứng dụng web theo kiến trúc client-server,
  sử dụng bộ công nghệ MERN Stack gồm MongoDB, Express.js, React.js và Node.js.
  Lý do lựa chọn hướng tiếp cận này xuất phát từ ba yếu tố chính. Thứ nhất, toàn
  bộ bộ công nghệ đều sử dụng JavaScript, giúp giảm chi phí chuyển đổi ngữ cảnh
  giữa phát triển giao diện và máy chủ. Thứ hai, React.js cho phép xây dựng giao
  diện dạng Single Page Application (SPA) với trải nghiệm người dùng mượt mà và
  khả năng quản lý trạng thái tập trung thông qua Redux. Thứ ba, MongoDB với mô
  hình tài liệu linh hoạt phù hợp với cấu trúc dữ liệu sản phẩm có nhiều biến thể
  và cấu hình vận chuyển phân cấp địa lý.

  Giải pháp của đồ án là ứng dụng web, bao gồm hai thành phần chính.
  Phần giao diện người dùng được phát triển bằng React.js kết hợp Redux, cho phép
  khách hàng duyệt sản phẩm, quản lý giỏ hàng, đặt hàng và theo dõi lịch sử giao
  dịch. Phần máy chủ được phát triển bằng Node.js và Express.js, cung cấp RESTful
  API xác thực bằng JWT, tích hợp cổng thanh toán VNPay, giao tiếp thời gian thực
  qua Socket.IO và gửi email thông báo qua Gmail OAuth2.

  Đóng góp chính của đồ án là một hệ thống thương mại điện tử hoàn chỉnh với
  nghiệp vụ thực tế, trong đó nổi bật là ba đóng góp đặc thù: quy trình hoàn trả
  hàng có vòng đời trạng thái đầy đủ từ khi khách yêu cầu đến khi hoàn tiền; hệ
  thống phân quyền theo vai trò kết hợp nhật ký kiểm toán hành động giúp truy vết
  mọi thao tác của nhân viên; và cấu hình phí vận chuyển linh hoạt đến cấp xã/phường.
  Kết quả đạt được là một ứng dụng web có thể triển khai thực tế, đáp ứng đầy đủ
  nghiệp vụ vận hành của một cửa hàng thời trang trực tuyến quy mô vừa và nhỏ.

  \section{Bố cục đồ án}
  \label{section:1.4}
  Phần còn lại của báo cáo đồ án tốt nghiệp này được tổ chức như sau.

  Chương 2 trình bày quá trình khảo sát hiện trạng và phân tích yêu cầu hệ thống.
  Từ kết quả khảo sát, chương này xác định các tác nhân và ca sử dụng của hệ thống
  thông qua biểu đồ ca sử dụng tổng quát và các biểu đồ usecase chi tiết cho
  từng nhóm chức năng. Tiếp đến, chương đặc tả các nghiệp vụ chính của hệ thống
  bao gồm quy trình đặt hàng, thanh toán, hoàn trả hàng và phân quyền nhân viên,
  và kết thúc bằng các yêu cầu phi chức năng về hiệu năng, bảo mật và khả năng
  mở rộng.

  Trong Chương 3, đồ án trình bày tổng quan về các công nghệ và thư viện được sử
  dụng trong quá trình xây dựng hệ thống. Chương này phân tích lý do lựa chọn từng
  thành phần trong bộ công nghệ MERN Stack, đồng thời giới thiệu các thư viện bổ
  trợ như Redux, Socket.IO, JWT và VNPay SDK, làm cơ sở lý thuyết cho các quyết
  định thiết kế được trình bày ở các chương sau.

  Chương 4 trình bày toàn bộ quá trình thiết kế và hiện thực hệ thống. Chương này
  bao gồm thiết kế kiến trúc tổng thể theo mô hình client-server, thiết kế cơ sở
  dữ liệu với các thực thể và mối quan hệ chính, thiết kế giao diện người dùng cho
  cả hai vai trò khách hàng và quản trị viên, và các kết quả kiểm thử chức năng
  đối với các ca sử dụng quan trọng.

  Chương 5 trình bày các giải pháp đóng góp của đồ án. Chương này bao gồm những
  giải pháp đưa ra giúp giải quyết được các bài toán giúp cải thiện những bài toán
  thực tiễn có trong đề tài.

  Chương 6 trình bày kết luận của đồ án, tổng hợp các kết quả đạt được so với mục
  tiêu đề ra ban đầu, đánh giá những hạn chế còn tồn tại và đề xuất hướng phát
  triển trong tương lai.
  
\end{document}